\section{Tinjauan Pustaka}

DeepBach \parencite{hadjeres2017} adalah model generatif musik simbolik berbasis arsitektur LSTM terarah yang dirancang khusus untuk menghasilkan \textit{four-part chorale} SATB dalam gaya Bach. Dipresentasikan di ICML 2017, model ini memperkenalkan paradigma generasi musik terkelola (\textit{steerable generation}) di mana pengguna dapat menetapkan batasan pada suara atau posisi tertentu. Penelitian ini menempatkan DeepBach sebagai titik komparatif historis awal, mewakili standar generasi \textit{chorale} simbolik berbasis data pada era \textit{deep learning} generasi pertama.

Coconet \parencite{huang2017} dikembangkan oleh Google Magenta. Model ini menggunakan pendekatan distribusi simultan seluruh suara (\textit{orderless NADE}) yang memungkinkan pelengkapan suara secara fleksibel dari kondisi parsial apa pun. Relevansi langsung Coconet terhadap penelitian ini terletak pada preseden metodologisnya: naskah asli Coconet menggunakan pustaka music21 untuk mengukur jumlah pelanggaran kuint sejajar dan oktaf sejajar pada \textit{output}-nya, membuktikan bahwa evaluasi berbasis kaidah harmoni formal merupakan metode yang valid dalam komunitas riset AI musik tingkat konferensi internasional.

NotaGen \parencite{wang2025} adalah model \textit{transformer} simbolik berskala besar yang dilatih pada lebih dari 1,6 juta karya musik dan \textit{fine-tuned} (disesuaikan lebih lanjut melalui pelatihan tambahan pada dataset spesifik) pada sekitar 9.000 komposisi klasik berkualitas tinggi. Diterima di IJCAI 2025, model ini mengklaim kemampuan generasi musik simbolik yang melampaui model sebelumnya dalam hal musikalitas dan koherensi partitur. Evaluasi dalam naskah aslinya mengandalkan penilaian estetika subjektif dan metrik keselarasan semantik (\textit{CLaMP 2 similarity score}), tanpa menyertakan analisis \textit{voice leading} berbasis kaidah harmoni formal. Celah ini menjadi salah satu kontribusi utama penelitian yang diusulkan.

Fang et al. \parencite{fang2020} dalam ``\textit{Bach or Mock? A Grading Function for Chorales in the Style of J.S. Bach}'' mengembangkan sistem penilaian algoritmis yang mengevaluasi \textit{output} model generatif terhadap kaidah harmoni Bach menggunakan pustaka music21. Penelitian ini merupakan preseden metodologis paling langsung bagi kerangka evaluasi yang dibangun dalam penelitian yang diusulkan: Fang et al. mendemonstrasikan bahwa model mutakhir, termasuk DeepBach sendiri, secara konsisten menghasilkan pelanggaran \textit{voice leading} yang terukur, bahkan ketika \textit{output} tersebut terdengar meyakinkan bagi pendengar non-ahli. Penelitian ini memformulasikan metrik kuantitatif yang mengukur rasio pelanggaran terhadap total momen harmonis, yang secara konseptual paralel dengan \textit{Strube Score} dalam penelitian yang diusulkan.

Le et al. \parencite{le2025} dalam survei komprehensif tentang metode \textit{natural language processing} untuk generasi dan \textit{information retrieval} musik simbolik mendokumentasikan pergeseran signifikan dalam komunitas MIR: dari evaluasi subjektif berbasis uji dengar (\textit{listening test}) menuju evaluasi algoritmis yang terukur dan dapat direplikasi. Survei ini menegaskan bahwa keterbatasan pendekatan subjektif, seperti variabilitas antar-pendengar dan rendahnya replikabilitas, mendorong kebutuhan akan kerangka evaluasi deterministik berbasis teori musik formal, yang menjadi dasar pendekatan penelitian yang diusulkan. Tren ini juga ditegaskan oleh temuan-temuan terbaru bahwa LLM musik masih memerlukan evaluasi pada aspek penempatan (\textit{reasoning}) spasial struktural dan kepatuhan (\textit{constraint adherence}) formal \parencite{Zhou2024, Ji2023}. Lebih jauh, perbandingan lintas arsitektur dalam konteks ini menekankan pentingnya isolasi faktor struktural melalui pendekatan berbasis kaidah dan pembobotan algoritmis \parencite{Chen2020}.

music21 \parencite{cuthbert2010} adalah pustaka Python sumber terbuka untuk musikologi berbasis komputer yang dikembangkan di MIT. Digunakan secara luas dalam komunitas MIR untuk analisis, manipulasi, dan evaluasi representasi musik simbolik. Penggunaan music21 sebagai instrumen evaluasi dalam penelitian ini mengikuti preseden yang telah ditetapkan oleh Huang et al. \parencite{huang2017} dan Fang et al. \parencite{fang2020}, dan memungkinkan verifikasi algoritmis yang konsisten terhadap seluruh 120 sampel MIDI \parencite{cuthbert2020}. Ketersediaan instrumen Python juga memungkinkan pengembangan yang lebih terarah \parencite{dong2020}.

Kebaruan Penelitian: Studi komparatif yang ada dalam literatur umumnya membandingkan model AI musik berdasarkan kualitas audio subjektif atau metrik statistik distribusional \parencite{Retkowski2024, wu2020}, bukan berdasarkan \textit{constraint adherence} terhadap kaidah harmoni teori musik formal mutlak. Berbeda dengan NotaGen yang pada evaluasi orisinalnya diandalkan murni pada kedekatan data tanpa memverifikasi pelanggaran \textit{rule-based} deterministik, penelitian ini menguji hipotesis bahwa absennya aturan negatif mutlak memicu kelalaian harmoni spasial yang sebelumnya dilaporkan pada riset LLM \parencite{Zhou2024, Ji2023}. Belum terdapat penelitian yang (1) membandingkan ketiga model ini secara bersamaan dalam satu kerangka evaluasi, (2) menggunakan teori Strube sebagai alat ukur yang tidak ambigu \parencite{choi2020}, dan (3) menguji NotaGen dalam konteks evaluasi harmoni formal deterministik absolut. Penelitian ini mengisi ketiga celah tersebut sekaligus \parencite{sturm2020}.

\section{Landasan Teori}

\subsection{Teori Harmoni Fungsional Gustav Strube}

Gustav Strube dalam \textit{The Theory and Use of Chords} \parencite{strube1928} menguraikan kaidah \textit{voice leading} yang membentuk fondasi harmoni fungsional Barat. Tiga kaidah yang dijadikan variabel dependen dalam penelitian ini adalah sebagai berikut.

\textit{Pertama}, larangan kuint sejajar (\textit{parallel fifths}): dua suara dilarang bergerak secara paralel dalam arah yang sama membentuk interval kuint sempurna di dua momen berurutan. Kaidah ini dijaga karena gerak kuint sejajar melemahkan independensi suara dan menciptakan kesan harmoni yang stagnan \parencite{harrison2020}.

\textit{Kedua}, larangan oktaf sejajar (\textit{parallel octaves}): dua suara dilarang bergerak secara paralel membentuk interval oktaf. Pelanggaran ini meniadakan kontribusi independen salah satu suara terhadap tekstur harmoni.

\textit{Ketiga}, kewajiban resolusi nada penuntun (\textit{leading tone resolution}): nada ketujuh dari tangga nada, yang memiliki tegangan harmonik tertinggi, harus bergerak ke atas menuju tonika (nada pertama) ketika frase harmonis mengindikasikan resolusi \parencite{schoenberg1954}. Kegagalan resolusi ini menghasilkan kesan harmoni yang menggantung atau tidak tuntas.

Ketiga kaidah ini dipilih karena sifatnya yang objektif dan tidak ambigu (masing-masing dapat diverifikasi secara biner pada setiap pasangan momen harmonis) serta posisinya sebagai bagian dari kurikulum resmi \programstudy{} \institutionname{}.

\subsection{Musik Simbolik dan Model Generatif Simbolik}

Musik simbolik merujuk pada representasi musik dalam bentuk notasi yang dapat dibaca secara struktural, seperti format MIDI atau \textit{ABC notation}. Representasi ini berbeda secara fundamental dari representasi audio (WAV, MP3): musik simbolik dapat diurai per suara, dianalisis struktur intervalnya, dan diverifikasi terhadap kaidah teori musik secara terprogram. Model generatif musik simbolik menghasilkan representasi ini, berbeda dari model audio populer seperti Suno yang menghasilkan berkas suara monolitik di mana pemisahan suara individual secara terprogram tidak dimungkinkan. Perbedaan paradigmatik inilah yang menjadi dasar fundamental mengapa penelitian ini hanya mengikutsertakan model simbolik.

\subsection{Evaluasi Otomatis dalam \textit{Music Information Retrieval} (MIR)}

\textit{Music Information Retrieval} adalah bidang interdisipliner yang mengembangkan metode komputasional untuk menganalisis, mengekstrak, dan mengklasifikasi informasi dari musik. Evaluasi berbasis kaidah harmoni formal dalam MIR memungkinkan analisis yang konsisten, terulang, dan berskala besar (dalam hal ini 120 sampel MIDI) tanpa bergantung pada subjektivitas pendengar sebagai sumber variabilitas. Pendekatan ini telah digunakan dalam penelitian fondasional seperti Coconet \parencite{huang2017} dan \textit{Bach or Mock} \parencite{fang2020}, dan menjadi standar yang semakin diakui dalam komunitas MIR internasional sebagai respons terhadap keterbatasan inheren evaluasi perseptual subjektif \parencite{yang2020, le2025}.

\subsection{Formalisasi Matematis Ruang Kosakata dan Arsitektur Model}

Sebuah sekuens musik simbolik $\mathbf{X}$ sepanjang $T$ dibentuk dari token diskrit yang merepresentasikan atribut musikal. Ruang kosakata multi-atribut $\mathbf{V}$ difaktorkan sebagai:

\[
\mathbf{V} = \mathbf{P} \times \mathbf{D} \times \mathbf{V}_{\text{idx}} \times \mathbf{M}
\]

Di mana $\mathbf{P} = \{0, 1, \dots, 127\}$ adalah ruang \textit{pitch} berdasarkan MIDI \textit{note number}, $\mathbf{D}$ adalah ruang durasi diskrit, $\mathbf{V}_{\text{idx}} = \{0,1,2,3\}$ adalah indeks suara SATB, dan $\mathbf{M}$ adalah posisi metrikal dalam takar.

Ketiga arsitektur yang dibandingkan menyelesaikan masalah dependensi dua poros (horizontal melodi dan vertikal harmoni) dengan pendekatan matematis yang berbeda.

\textit{Recurrent} (DeepBach) mengestimasi distribusi probabilitas kondisional lokal per koordinat menggunakan LSTM dua arah:

\[
P(x_{v,t} \mid \mathbf{X}_{\setminus \{v,t\}}, \mathbf{M}) = \sigma\left(\Phi(h_{v,t}^{\rightarrow}, h_{v,t}^{\leftarrow}, \Omega(\mathbf{X}_{\text{concurrent}}))\right)
\]

Konvolusional (Coconet) memperlakukan partitur sebagai matriks spasial dan melakukan \textit{infilling} melalui \textit{Orderless Neural Autoregressive Density Estimator} pada wilayah bermasker $\mathbf{M}_{\text{mask}}$:

\[
\hat{\mathbf{X}} = f_{\theta}(\mathbf{X} \odot \mathbf{M}_{\text{mask}} + \mathbf{M}_{\text{mask}})
\]

\textit{Attention-based} (NotaGen) menserialisasikan partitur menjadi teks satu dimensi dan mengandalkan \textit{global self-attention} (mekanisme perhatian yang menelusuri ketergantungan antar seluruh elemen dalam satu sekuens sekaligus) untuk menangkap dependensi jarak jauh, baik vertikal maupun horizontal:

\[
\text{Attention}(Q, K, V) = \text{softmax}\left(\frac{QK^T}{\sqrt{d_k}}\right)V
\]

Perbedaan pendekatan matematis ini menjadi dasar hipotesis bahwa \textit{bias} induktif arsitektur merupakan faktor utama yang memengaruhi pola \textit{constraint adherence} pada \textit{output} masing-masing model. Perlu diakui bahwa ketiga model juga berbeda dalam skala data pelatihan dan spesialisasi tugas: DeepBach dan Coconet dilatih khusus pada korpus \textit{chorale} SATB Bach ($\sim$400 karya), sementara NotaGen dilatih pada 1,6 juta partitur lintas genre dan dikondisikan ke tugas SATB melalui metadata teks. Heterogenitas ini menjadi \textit{acknowledged confound} dalam desain komparatif penelitian ini: perbedaan skor antar model kemungkinan mencerminkan kombinasi antara \textit{bias} induktif arsitektur dan skala/spesialisasi data pelatihan, bukan arsitektur semata. Interpretasi temuan akan mempertimbangkan dua penjelasan alternatif ini secara eksplisit.
